\documentclass[a4paper,12pt]{article}
\usepackage[utf8]{inputenc}
\usepackage{amsmath}
\usepackage{graphicx}
\usepackage{float}
\usepackage{hyperref}
\usepackage{amsfonts}
\usepackage{amssymb}
\usepackage{geometry}
\usepackage{listings}
\usepackage{xcolor}
% Pacote para códigos em Python
\usepackage{listings}
\renewcommand{\lstlistingname}{Algoritmo}
% Configuração do ambiente listings para Python
\lstset{
    language=Python,
    basicstyle=\ttfamily\small,
    keywordstyle=\color{blue}\bfseries,
    stringstyle=\color{red},
    commentstyle=\color{gray}\itshape,
    numbers=left,
    numberstyle=\tiny\color{gray},
    stepnumber=1,
    numbersep=5pt,
    backgroundcolor=\color{white},
    frame=single,
    tabsize=4,
    showspaces=false,
    showstringspaces=false,
    breaklines=true,
    breakatwhitespace=true,
    captionpos=b
}
\pagestyle{empty}

\geometry{a4paper, margin=0.5in,top=0.2in}

\title{Prova 2 - Álgebra Linear para Computação}
\date{}
\newcommand{\avg}{\mathbf{\textbf{avg}}}
\newcommand{\std}{\mathbf{\textbf{std}}}
\begin{document}
\maketitle
\vspace{-4em}
\begin{enumerate}
    \item \texttt{(2 pontos)} Considere uma matriz $A \in \mathbb{R}^{2\times 2}$. Aplique a decomposição $LU$ e determine as 4 equações envolvendo os elementos $\ell_{ij} \in L$ e $u_{ij} \in U$. Após isso, faça a decomposição $LU$ e utilize a decomposição para calcular o determinante de $A$.
    $$
    A = \begin{bmatrix}
        5 & 2 \\
        3 & 1  \\
    \end{bmatrix}.
    $$

    \item \texttt{(2 pontos)} Utilizando o terorema espectral:
    
    \begin{enumerate}
        \item \texttt{(1.5 ponto)} faça os cálculos passo a passo para a decomposição $A = P D P^T$ onde:
    
        $$
        \underbrace{\begin{bmatrix}
            6 & 2\\
            2 & 3\\
        \end{bmatrix}}_{A} =  \underbrace{\frac{1}{\sqrt{5}}\begin{bmatrix}
            2 & 1\\
            1 & -2\\
        \end{bmatrix}}_{P}\ \underbrace{\begin{bmatrix}
            7 & 0\\
            0 & 2\\
        \end{bmatrix}}_{D} \ \underbrace{\frac{1}{\sqrt{5}}\begin{bmatrix}
            2 & 1\\
            1 & -2\\
        \end{bmatrix}}_{P^T}.
        $$
        \item \texttt{(0.25 ponto)} Oquê acontece quando calculamos $P P^T$?
        \item \texttt{(0.25 ponto)} Utilize o resultado acima para calcular o determinante e o trço de $A$.
    \end{enumerate}
    

    \item \texttt{(2 pontos)} Dada a autodecomposição $A = P D P^{-1}$, explique as caracteristicas das matrizes $P$ e $D$. Depois, calcular a autodecomposição da matriz:
    $$
    A = \begin{bmatrix}
        1 & 4\\
        2 & 3\\
    \end{bmatrix}.
    $$
    Calcule $A^4$ e explique como relaizar os cálculos utilizando o Teorema da Autodecomposição.

    \item \texttt{(2 pontos)} No caso especial $n = 1$, o problema geral de mínimos quadrados reduz-se a encontrar um escalar $x$ que minimize $\|\mathbf{a}x - \mathbf{b}\|^2$, onde $\mathbf{a}$ e $\mathbf{b}$ são vetores em $\mathbb{R}^m$. (Escrevemos a matriz $A$ em letra minúscula aqui, pois trata-se de um vetor de dimensão $m$.) Assumindo que $\mathbf{a}$ e $\mathbf{b}$ são não nulos, mostre que
    $$
    \|\mathbf{a}x - \mathbf{b}\|^2 = \|\mathbf{b}\|^2 (\sin \theta)^2,
    $$
    onde $\theta = \angle(\mathbf{a}, \mathbf{b})$.
    

    \item \texttt{(2 pontos)} No problema de \emph{mínimos quadrados ponderados}, minimizamos o objetivo
    \[
    \sum_{i=1}^m w_i (\tilde{a}_i^T x - b_i)^2,
    \]
    onde $w_i$ são pesos positivos dados.  
    Os pesos permitem atribuir diferentes importâncias às componentes do vetor de resíduos.  
    (O objetivo do problema ponderado é o quadrado da norma ponderada,
    $\|Ax - b\|_w^2$)

    \begin{enumerate}
        \item \texttt{(1 ponto)} Mostre que o objetivo de mínimos quadrados ponderados pode ser expresso como
        \[
        \|D(Ax - b)\|^2,
        \]
        para uma matriz diagonal apropriada $D$.  
        Isso nos permite resolver o problema ponderado como um problema padrão de mínimos quadrados,  
        minimizando $\|Bx - d\|^2$, onde $B = DA$ e $d = Db$.

        \item \texttt{(1 ponto)} A solução aproximada de mínimos quadrados é dada por  
        \[
            \hat{x} = (A^T A)^{-1} A^T b.
        \]
        Dê uma fórmula semelhante para a solução do problema de mínimos quadrados ponderados. Você pode utilizar $W = \text{diag}(w)$
    \end{enumerate}
\end{enumerate}

\end{document}